\section{Science Program Observations}
\label{sec:science_program}
%
The science program observations acquired during Rubin commissioning, which form the basis of \gls{DP2}, spanned three distinct phases.
The first comprised small-field observations acquired shortly after the \gls{LSSTCam} First Photon milestone in support of the Rubin First Look event (\ref{ssec:small_fields}).
The second was the commissioning-era \gls{SV} survey campaigns, which ran from \svcampaignstartdate to \svcampaignenddate (\ref{ssec:sv_surveys}).
The end of the \gls{SV} campaign was followed by a four-week engineering downtime, after which the observatory was formally handed over from the Construction project to Rubin Operations on \startoperations.
The third phase comprises pre-\gls{LSST} Early Operations observations from that date through the \gls{DP2} cutoff date of \dptwoenddate.
Early Operations observations continued for approximately six months after the \gls{DP2} cutoff, until the start of the \gls{LSST} on \lsststartdate.
All science program observations were scheduled by the \gls{FBS}, which selected visits adaptively to balance sky coverage, filter distribution, airmass, and lunar phase constraints \citep[][]{2026SPIE14151E..0XJ}.
Because the program was conducted while the system was still being commissioned and optimized, visit acquisition was driven by engineering and commissioning priorities rather than by a scientifically optimized observing strategy.
Engineering observations are not included in \gls{DP2}.

\figref{fig:dp2_mjd_cumulative} shows the cumulative distribution of visit MJD for all \gls{DP2} visits, combined across all filters.
Notable dates are indicated: the \gls{DP2} start date, \dptwostartdate; the end of the \gls{DP2} data collection window, \dptwoenddate; a phase of rapid data acquisition (July 1--25, 2025); and the end of the engineering shutdown, when observations resumed at the start of Operations, \startoperations.
The sky distribution of \gls{DP2} visits is shown in \figref{fig:dp2_sky}.
The observation cadence during \gls{DP2} differs significantly from that of the planned \gls{LSST}.
Rather than following the cadence intended for the full survey, visit acquisition during this period was driven by engineering and commissioning priorities.
As a result, the distribution of visits in time reflects the sequence of commissioning activities and engineering interventions rather than a scientifically-optimized observing strategy.
\begin{figure}[hbt!]
\plotone{dp2_visit_mjd_cumulative}
\caption{Cumulative distribution of visit MJD for all \gls{DP2} visits, combined across all filters. Vertical dotted lines mark notable dates: the \gls{DP2} start date, \dptwostartdate; the end of the data collection window, \dptwoenddate; a phase of rapid data acquisition (July 1--25, 2025); and the end of the engineering shutdown, when observations resumed \startoperations}
\label{fig:dp2_mjd_cumulative}
\end{figure}

\begin{figure*}[hbt!]
\plotone{dp2_sky_coverage}
\caption{Sky distribution of \gls{DP2} visits across the three components of the observing program. Small-field survey areas are marked with black circles and \glspl{DDF} with violet squares. The \gls{SV} wide survey covers \coaddedarea tracing the ecliptic plane, with visit density increasing toward the 300\,deg$^2$ core. Pre-\gls{LSST} Early Operations observations appear as sparser, predominantly $i$-band coverage, extending across a larger area of the \gls{LSST} footprint.}
\label{fig:dp2_sky}
\end{figure*}

Data collected between the end of the \gls{DP2} observing window on \dptwoenddate and the formal start of the \gls{LSST} on \lsststartdate will be processed and made available to \gls{LSST} data rights holders through the same \gls{RSP} services.
The format and timing of that release are still to be determined; planning will begin once \gls{DP2} has been delivered \citep[][]{RTN-011}

% Small fields
\subsection{Small-Field Observations}
\label{ssec:small_fields}
%
Several small fields were observed prior to the start of the \gls{SV} surveys, with some observations overlapping the early portion of the \gls{SV} survey period.
All small-field observations, including those from the First Look campaign, are included in \gls{DP2}.
Observations were typically acquired in sequences of at least 10 visits per bandpass, cycling through the available filters, with field selection guided by commissioning priorities and community input.
Several of these fields, including \gls{COSMOS}, M49, the Prawn Nebula region, and Rubin\_SV\_225\_-40, received substantial numbers of visits and are among the most deeply observed regions in \gls{DP2}.
The \gls{COSMOS} field was first observed as part of the small-field campaign but was not accessible during the main \gls{SV} survey period; additional observations were later acquired during the Early Operations phase.

\tabref{tab:dp2_sf_m5_depths} reports the 5$\sigma$ point--source depths for coadded images per field and per band, where coverage in a band is nonzero.
% together with the expected 10 yr LSST depths derived from the baseline simulated survey \citep[][]{2022ApJS..258....1B}.

%%%%% This table is auto generated from data, DO NOT EDIT
\setlength{\tabcolsep}{3pt}
\begin{deluxetable}{lcccccc}
\caption{Median $5\sigma$ coadd detection limits (AB mag) per band per small-field science verification region in DP2, estimated as the median PSF magnitude of point sources (extendedness $\leq 0.5$) with $4.9 < {\rm S/N} < 5.1$.}
\label{tab:dp2_sf_m5_depths}
\tablehead{
  \colhead{\textbf{Field Code}} & \multicolumn{6}{c}{\textbf{Band}}\\
  \cline{2-7}
   & $u$ & $g$ & $r$ & $i$ & $z$ & $y$ 
}
\startdata
Abell\_2764 & \nodata & 25.80 & \nodata & 22.72 & \nodata & \nodata \\
DESI\_SV3\_R1 & \nodata & \nodata & \nodata & \nodata & \nodata & \nodata \\
M49 & 25.57 & 26.87 & 26.59 & 25.95 & \nodata & \nodata \\
Prawn & 25.14 & 26.07 & 25.82 & 25.26 & 23.87 & \nodata \\
Trifid-Lagoon & 26.09 & 26.76 & 25.96 & 25.34 & 23.61 & 22.34 \\
New\_Horizons & 25.21 & 26.29 & 26.33 & 25.93 & 25.14 & 23.61 \\
Rubin\_SV\_212\_-7 & \nodata & 26.95 & 26.72 & 26.25 & \nodata & \nodata \\
Rubin\_SV\_216\_-17 & \nodata & 26.32 & 26.08 & 25.93 & \nodata & \nodata \\
Rubin\_SV\_225\_-40 & 26.06 & 27.09 & 26.53 & 26.16 & 25.34 & 23.89 \\
Rubin\_SV\_280\_-48 & 23.79 & 24.80 & 24.67 & 24.67 & 23.53 & \nodata \\
Rubin\_SV\_300\_-41 & \nodata & 24.59 & \nodata & \nodata & 23.65 & \nodata \\
Rubin\_SV\_320\_-15 & 24.24 & 25.06 & 24.90 & 25.07 & 24.31 & 22.68 \\
\enddata
\end{deluxetable}


%%%%%  SV commissioning surveys
\subsection{Science Validation Surveys}
\label{ssec:sv_surveys}
%
The \gls{SV} surveys began on \svcampaignstartdate and comprised two interleaved components within a single \gls{FBS} configuration, a wide survey and a deep survey.

\subsubsection{The Wide Survey}
The Wide Survey initially covered a $\sim$\,3,000\,deg$^2$ footprint tracing the ecliptic plane from the dense regions of the Galactic Bulge through low-dust regions of the planned \gls{LSST} \gls{WFD} footprint.
Approximately halfway through the \gls{SV} period, on \svwfdsevenfiftystart{}, the footprint was consolidated to $\sim$\,750\,deg$^2$ and observed in the $griz$ bands only.
This reduced footprint is a $10\degr$-wide band centred on the ecliptic, defined by ecliptic longitude $285\degr < \lambda < 360\degr$ and ecliptic latitude $|\beta| < 5\degr$.
On \svwfdthreehundredstart{}, in the final week of commissioning, the footprint was further reduced to the $\sim$\,300\,deg$^2$ area covered by newly deployed templates ($300\degr <$ R.A. $< 324\degr$, $-26\degr <$ decl. $< -10\degr$, and $\beta < 5\degr$), to concentrate visits and reach the depth required for difference imaging tests.
As a result, median visit counts and coadded depth increase substantially toward the inner regions of the footprint.

\tabref{tab:dp2_sv_wide_footprints} provides a description of the WFD areas observed during the campaign.

%%%%% This table is auto generated from data, DO NOT EDIT
\setlength{\tabcolsep}{3pt}
\begin{deluxetable*}{lllc}
\tablewidth{0pt}
\caption{Definitions of the successive footprints of the \gls{SV} Wide Survey. $\lambda$ and $\beta$ are ecliptic longitude and latitude in the barycentric true ecliptic frame. Areas are computed on a HEALPix grid with $N_{\rm side}=256$. About 77\% of the 300\,deg$^2$ footprint lies within the 750\,deg$^2$ footprint and 99\% within the 3,000\,deg$^2$ footprint.}
\label{tab:dp2_sv_wide_footprints}
\tablehead{
  \colhead{\textbf{Footprint}} & \colhead{\textbf{Definition}} & \colhead{\textbf{Extent}} & \colhead{\textbf{Area (deg$^2$)}}\\
}
\startdata
3,000\,deg$^2$ & $|\beta| < 10\degr$; $\lambda > 240\degr$ or $\lambda < 40\degr$ & $20\degr$ wide in $\beta$, $160\degr$ long in $\lambda$ & 3,184 \\
750\,deg$^2$ & $|\beta| < 5\degr$; $\lambda > 285\degr$ & $10\degr$ wide in $\beta$, $75\degr$ long in $\lambda$ & 749 \\
300\,deg$^2$ & $300\degr <$ R.A. $< 324\degr$; $-26\degr <$ decl. $< -10\degr$; $\beta < 5\degr$ & $\lambda \approx 297\degr$--$323\degr$, $\beta \approx -11\degr$ to $+5\degr$ & 305 \\
\enddata
\end{deluxetable*}


\tabref{tab:dp2_sv_wide_m5_depths} provides the estimated $m_5$ coadded depths within each of the subsets of the wide SV area.

%%%%% This table is auto generated from data, DO NOT EDIT
\setlength{\tabcolsep}{3pt}
\begin{deluxetable}{lcccccc}
\caption{Median $5\sigma$ coadd detection limits (AB mag) per band in disjoint areas of the wide-area \gls{SV} survey footprint: the 300\,deg$^2$ core, the 750\,deg$^2$ area outside it, and the 3,000\,deg$^2$ area outside both in DP2, estimated as the median PSF magnitude of point sources (extendedness $\leq 0.5$) with $4.9 < {\rm S/N} < 5.1$.}
\label{tab:DP2_sv_wide_m5_depths}
\tablehead{
  \colhead{\textbf{Field Code}} & \multicolumn{6}{c}{\textbf{Band}}\\
  \cline{2-7}
   & $u$ & $g$ & $r$ & $i$ & $z$ & $y$ 
}
\startdata
3,000\,deg$^2$ excl.\ 750 & 24.16 & 25.21 & 24.73 & 24.32 & 23.68 & 22.47 \\
750\,deg$^2$ excl.\ 300 & 24.11 & 25.32 & 25.07 & 24.68 & 23.85 & 22.74 \\
300\,deg$^2$ & 24.27 & 25.00 & 24.84 & 24.75 & 24.03 & 22.85 \\
\enddata
\end{deluxetable}


%%%% Deep survey
\subsubsection{The Deep Survey}
The Deep Survey targeted four of the five planned \gls{LSST} \gls{DDF}s: \gls{ECDFS}, \gls{EDFS} (fields A and B), ELAISS1, and \gls{COSMOS}.
The \gls{FBS} interleaved \gls{DDF} observations with the wide survey, scheduling visits to each field according to the prescribed \gls{DDF} observing cadence.
ELAISS1 received the largest number of visits, 539 across all six $ugrizy$ bands, and is the only \gls{DDF} observed in all bands during the \gls{SV} survey.
\gls{COSMOS} was not accessible during the \gls{SV} period as described above.
The remaining \glspl{DDF} were observed primarily in $griz$, as the $u$ and $y$ filters were unavailable after 2025-08-10 due to a hardware issue with the filter exchanger.
The dither pattern for the \glspl{DDF} evolved during the campaign; early observations used only small translational dithers with no intra-night rotation, while later observations introduced both translational and rotational intra-night dithers.

\tabref{tab:dp2_ddf_m5_depths} reports the 5$\sigma$ point--source depths for coadded images per field and per band, where coverage in a band is nonzero.

%%%%% This table is auto generated from data, DO NOT EDIT
\setlength{\tabcolsep}{3pt}
\begin{deluxetable}{lcccccc}
\caption{Median $5\sigma$ coadd detection limits (AB mag) per band per Deep Drilling Field in DP2, estimated as the median PSF magnitude of point sources (extendedness $\leq 0.5$) with $4.9 < {\rm S/N} < 5.1$.}
\label{tab:dp2_ddf_m5_depths}
\tablehead{
  \colhead{\textbf{Field Code}} & \multicolumn{6}{c}{\textbf{Band}}\\
  \cline{2-7}
   & $u$ & $g$ & $r$ & $i$ & $z$ & $y$ 
}
\startdata
ELAIS\_S1 & 25.35 & 26.67 & 26.04 & 26.22 & 25.27 & 23.81 \\
XMM\_LSS & \nodata & \nodata & \nodata & \nodata & \nodata & \nodata \\
ECDFS & 24.39 & 25.96 & 25.67 & 26.00 & 24.92 & 22.51 \\
EDFS\_a & 23.64 & 25.75 & 25.61 & 25.79 & 24.71 & 23.22 \\
EDFS\_b & 23.71 & 25.86 & 25.62 & 25.80 & 24.74 & 23.19 \\
COSMOS & 25.87 & 26.84 & 26.51 & 26.09 & 25.15 & 23.81 \\
\enddata
\end{deluxetable}


% Early ops
% TODO Say something about the non-coadded area -- e.g median five sigma depth
\subsection{Early Operations Observations}
\label{ssec:early_ops}
%
Following the transition to Operations on \startoperations, \gls{FBS}-driven observations resumed during the Early Operations optimization period.
Observations during this period, which ran from \startoperations through \dptwoenddate, were a mixture of engineering activities and \gls{FBS}-driven survey observations aimed at improving image quality and survey speed in preparation for the start of the \gls{LSST}, and building difference-imaging templates across the full \gls{LSST} footprint.

The pre-\gls{LSST} observations are characterised by a lower and less regular cadence than the \gls{SV} surveys, reflecting frequent engineering interruptions.
The \gls{FBS} was configured with a single-visit template tier as a scheduling priority, and on any given night would typically concentrate observations in a single band rather than interleaving filters, to maximise sky coverage for template building.
The $i$ band was preferred due to its high throughput and suitability for difference imaging.
A concentrated block of $i$-band observations in late October 2025 reflected deliberate \gls{AOS} engineering tests that held the filter fixed to isolate optical state changes from filter-induced perturbations, producing the prominent $i$-band coverage regions visible in \figref{fig:dp2_sky}.
Pre-\gls{LSST} visits not overlapping existing \gls{SV} survey regions were not combined into coadds but are included in the \gls{DP2} visit dataset.

% Image Quality
\subsection{Delivered Image Quality}
\label{ssec:image_quality}
%
The delivered image quality at Rubin Observatory reflects contributions from the atmosphere, dome environment, telescope optics, \gls{AOS}, and LSSTCam.
The design specification for the median delivered \gls{PSF} \gls{FWHM}, referenced to zenith at 550\,nm, is 0\farcs7 \citep[][]{LPM-17}, corresponding to the measured median atmospheric seeing of 0\farcs72 at Cerro Pach\'on \citep[][]{SOTN-004} added in quadrature with a total system contribution of 0\farcs35 \citep[][]{SOTN-004}.
During the period covered by the \gls{DP2} dataset, the system had not yet consistently met this design specification \citep[][]{SOTN-004, SITCOMTN-170}.
The primary excess contributions are from the optical system, including residual \gls{AOS} errors and mirror thermal gradients, as well as from dome seeing.
The current total system contribution to \gls{PSF} \gls{FWHM} is approximately 0\farcs55--0\farcs60, compared to the design allocation of 0\farcs40.
At the median survey airmass of 1.2, an atmospheric contribution of 0\farcs72 at zenith corresponds to $\sim$\,0\farcs80; combined in quadrature with the current system contribution, this is consistent with the observed median delivered image quality of $\sim$\,1\farcs0 \gls{PSF} \gls{FWHM} in the $riz$ bands during the \gls{DP2} observations \citep[][]{RTN-122}.

\figref{fig:delivered_image_quality} shows the cumulative distribution of the median PSF FWHM (in arcseconds) per visit computed across all detectors, for each band.
\begin{figure}[hbt!]
\plotone{dp2_median_per_visit_psf_fwhm_ecdf}
\caption{Cumulative distribution of the median per-visit seeing (PSF FWHM) per band. All regions included.}
\label{fig:delivered_image_quality}
\end{figure}
